
% --- Fix Horizontal Rules (---) ---
% Save the original \rule command so we can still use it later
\let\oldrule\rule

% Redefine \rule to check if it's the specific one generated by Pandoc for '---'
\renewcommand{\rule}[2]{%
  % Check: Is the requested width exactly 50% of the line width?
  % (Pandoc converts markdown '---' to \rule{0.5\linewidth} by default)
  \ifdim#1=0.5\linewidth
    \hrule % If yes, replace it with a full-width horizontal rule
  \else
    \oldrule{#1}{#2} % If no, draw the rule exactly as requested originally
  \fi
}

% --- Force Equation Numbering ---
% Redefine standard display math \[ ... \] to use the numbered 'equation' environment
% This ensures $$...$$ in Markdown gets numbered in the PDF
\renewcommand{\[}{\begin{equation}}
\renewcommand{\]}{\end{equation}}

\usepackage[most]{tcolorbox}

% --- Left-align all captions ---
% singlelinecheck=false stops LaTeX from centring captions that happen to fit
% on one line; justification=raggedright left-aligns the rest.
\usepackage[justification=raggedright,singlelinecheck=false]{caption}

% --- Left-align all section headings ---
% The KOMA classes that Quarto books use set headings with \raggedsection;
% make it explicitly ragged-right, so that a wrapped heading is never
% justified or centred.
\makeatletter
\@ifundefined{raggedsection}{}{\renewcommand*{\raggedsection}{\raggedright}}
\makeatother


% --- Color Definitions ---
% Theorem group (Theorem, Lemma, Corollary, Proposition)
\definecolor{thm-bg}{HTML}{F0F7FF}
\definecolor{thm-frame}{HTML}{CFE2FF}

% Definition
\definecolor{def-bg}{HTML}{F0FFF4}
\definecolor{def-frame}{HTML}{BADBCC}

% Example
\definecolor{ex-bg}{HTML}{FFFBF0}
\definecolor{ex-frame}{HTML}{FFEEBB}

% Proof
\definecolor{proof-bg}{HTML}{E2E3E5}
\definecolor{proof-frame}{HTML}{E2E3E5}

% Solution
\definecolor{sol-bg}{HTML}{F9F2F4}
\definecolor{sol-frame}{HTML}{D9534F}
% Algorithm
\definecolor{alg-bg}{HTML}{F3E5F5}    % Light Purple background
\definecolor{alg-frame}{HTML}{CE93D8} % Darker Purple frame

% --- Box Environments ---

% Theorem
\tcolorboxenvironment{theorem}{
  colback=thm-bg,
  colframe=thm-frame,
  boxrule=0pt,
  sharp corners,
  left=5pt, right=5pt,
  top=5pt, bottom=5pt,
  breakable
}

% Lemma
\tcolorboxenvironment{lemma}{
  colback=thm-bg,
  colframe=thm-frame,
  boxrule=0pt,
  sharp corners,
  left=5pt, right=5pt,
  top=5pt, bottom=5pt,
  breakable
}

% Corollary
\tcolorboxenvironment{corollary}{
  colback=thm-bg,
  colframe=thm-frame,
  boxrule=0pt,
  sharp corners,
  left=5pt, right=5pt,
  top=5pt, bottom=5pt,
  breakable
}

% Proposition
\tcolorboxenvironment{proposition}{
  colback=thm-bg,
  colframe=thm-frame,
  boxrule=0pt,
  sharp corners,
  left=5pt, right=5pt,
  top=5pt, bottom=5pt,
  breakable
}

% Definition
\tcolorboxenvironment{definition}{
  colback=def-bg,
  colframe=def-frame,
  boxrule=0pt,
  sharp corners,
  left=5pt, right=5pt,
  top=5pt, bottom=5pt,
  breakable
}

% Example
\tcolorboxenvironment{example}{
  colback=ex-bg,
  colframe=ex-frame,
  boxrule=0pt,
  sharp corners,
  left=5pt, right=5pt,
  top=5pt, bottom=5pt,
  breakable
}

% Proof
\tcolorboxenvironment{proof}{
  colback=proof-bg,
  colframe=proof-frame,
  boxrule=0pt,
  sharp corners,
  left=5pt, right=5pt,
  top=5pt, bottom=5pt,
  breakable
}

% Solution
\tcolorboxenvironment{sol}{
  colback=sol-bg,
  colframe=sol-frame,
  boxrule=0pt,
  sharp corners,
  left=5pt, right=5pt,
  top=5pt, bottom=5pt,
  breakable
}

% 2. Define the base environment as a simple "block" (non-floating)
%\newenvironment{algorithm}{}{}
% 3. Apply the tcolorbox styling (Header Banner Style)
% \tcolorboxenvironment{algorithm}{
%   colback=alg-bg,
%   colframe=alg-frame,
%   boxrule=0pt,
%   sharp corners,
%   left=5pt, right=5pt,
%   top=5pt, bottom=5pt,
%   breakable,
%   % --- Header Setup ---
%   fonttitle=\bfseries\sffamily\large, 
%   coltitle=black,
%   title=Algorithm,
%   colbacktitle=alg-frame, % Matches the banner to the frame color
%   titlerule=0pt,           % Removes the rule/line under the header text
%   toptitle=4pt,            % Adjusts padding at the top of the banner
%   bottomtitle=4pt          % Adjusts padding at the bottom of the banner
% }

\tcolorboxenvironment{algorithm}{
  colback=alg-bg,
  colframe=alg-frame,
  boxrule=0pt,
  sharp corners,
  left=5pt, right=5pt,
  top=5pt, bottom=5pt,
  breakable
}

% --- Code and printed output ---
% Wrap long source lines instead of running past the margin, and set code a
% size smaller and printed R output two sizes smaller, so that wide data frames
% (printed at width 110, see _quarto.yml) fit on one line.
% (\AtBeginDocument: Pandoc defines Highlighting after the header includes.)
\usepackage{fvextra}
\AtBeginDocument{%
  \RecustomVerbatimEnvironment{Highlighting}{Verbatim}{commandchars=\\\{\},breaklines,breakanywhere,fontsize=\small}%
  \RecustomVerbatimEnvironment{verbatim}{Verbatim}{breaklines,breakanywhere,fontsize=\scriptsize}%
}

% --- Fit gt tables to the text width ---
% gt sets its tables as \begin{tabular*}{width}{@{\extracolsep{\fill}}spec}, with
% fixed column widths that can add up to more than the line; gt-latex.lua
% rewrites them as \begin{gtfit}{width}{spec}. A table that fits is set at gt's
% width as before; a wider one is set at its natural width, and scaled down to
% the line width if needed.
\usepackage{environ}
\usepackage{adjustbox}
\newsavebox{\gtfitbox}
\NewEnviron{gtfit}[2]{%
  \sbox{\gtfitbox}{\begin{tabular}{#2}\BODY\end{tabular}}%
  \ifdim\wd\gtfitbox>#1\relax
    \adjustbox{max width=\linewidth}{\usebox{\gtfitbox}}%
  \else
    \begin{tabular*}{#1}{@{\extracolsep{\fill}}#2}\BODY\end{tabular*}%
  \fi
  \par}% end the line, so that gt's footnotes start below the table
